\subsection{Harmonic nuclear Hamiltonians}

On surface $s$, the harmonic Hamiltonian in mass-weighted normal coordinates
$Q_s$ is $H_s=E_s^{\min}+\tfrac12\sum_k(-\partial_{Q_{sk}}^2+\omega_{sk}^2 Q_{sk}^2)$
in atomic units. Here $E_s^{\min}$ is the electronic minimum, $Q_s$ is in
$\sqrt{m_e}\,\mathrm{bohr}$, and $\omega_{sk}$ denotes a vibrational energy
quantum in Hartree. For a PySCF molecule, the workflow uses its atomic
masses in atom order, including isotope assignments, unless explicit
masses are supplied. In the adiabatic Hessian (AH) model, the two Hessians
give separate frequencies and normal modes. After mass-weighted alignment
and removal of rigid motions, the orthogonal polar factor of the mode
overlap gives the Duschinsky matrix $J$ in $Q_f=JQ_i+K$. Projecting the
initial minimum relative to the aligned final minimum onto the final
modes gives $K$. The singular values of the mode overlap and the residual
of the displacement projection diagnose the affine approximation. The
signed vibronic origin is $\Delta_{00}=E_f^{\min}-E_i^{\min}+Z_f-Z_i$,
where $Z_s=\tfrac12\sum_k\omega_{sk}$.\cite{santoro}

The adiabatic-shift (AS) model keeps the displacement and energy difference
between the electronic minima and uses the ground-state curvature on both
surfaces. The vertical-gradient (VG) model also uses equal curvatures. It
places the excited-state minimum at $-\Omega_g^{-2}g$ in ground-state normal
coordinates, where $g$ is the projected excited-state gradient at the
ground-state minimum and $\Omega_g=\operatorname{diag}(\boldsymbol\omega_g)$.
Completing the square lowers the vertical gap by
$\tfrac12 g^{\mathsf T}\Omega_g^{-2}g$. The zero-point correction cancels in
both equal-curvature models.

For harmonic initial and final occupations $\mathbf n$ and $\mathbf m$, let
$F_{\mathbf n\mathbf m}=\langle\chi_{f\mathbf m}|\chi_{i\mathbf n}\rangle$.
With the combined index $\mathbf a=(\mathbf n,\mathbf m)$, the Hermite
generating functions of the two oscillator bases give
\begin{equation}
 G(z)=\sum_{\mathbf a}F_{\mathbf a}\frac{z^{\mathbf a}}{\sqrt{\mathbf a!}}
   =F_{\mathbf 0}\exp\bigl(\ell^{\mathsf T}z+\tfrac12 z^{\mathsf T}Bz\bigr).
 \label{eq:fc-generating}
\end{equation}
Writing $\Omega_s=\operatorname{diag}(\boldsymbol\omega_s)$, the product of
the two normalized vacuum Gaussians has precision
$P=\Omega_i+J^{\mathsf T}\Omega_f J$ and mean
$c=-P^{-1}J^{\mathsf T}\Omega_f K$ in the initial coordinates.
With $c_f=Jc+K$ and
$U=\bigl[(2\Omega_i)^{1/2};\,(2\Omega_f)^{1/2}J\bigr]$ denoting a vertical
stack, the Gaussian integral gives
\begin{equation}
 \begin{aligned}
 F_{\mathbf 0}&=
 \frac{2^{d/2}(\det\Omega_i\det\Omega_f)^{1/4}}{\sqrt{\det P}}
 \exp\!\left[-\tfrac12(c^{\mathsf T}\Omega_i c+c_f^{\mathsf T}\Omega_f c_f)\right],\\
 \ell&=\begin{pmatrix}(2\Omega_i)^{1/2}c\\(2\Omega_f)^{1/2}c_f\end{pmatrix},
 \qquad B=UP^{-1}U^{\mathsf T}-I ,
 \end{aligned}
\end{equation}
where $d$ is the number of modes on each surface. Differentiating
Eq.~\eqref{eq:fc-generating} gives, for any $a_j>0$,
\begin{equation}
 \sqrt{a_j}\,F_{\mathbf a}=
 \ell_j F_{\mathbf a-\mathbf e_j}
 +\sum_k B_{jk}\sqrt{a_k-\delta_{jk}}\,
 F_{\mathbf a-\mathbf e_j-\mathbf e_k} .
 \label{eq:fc-recurrence}
\end{equation}
The sum includes only occupied components of $\mathbf a-\mathbf e_j$.
Here $\mathbf e_j$ adds one quantum to component $j$. Thus a basis with a
total-quanta cutoff can be filled in order of increasing excitation.

\subsection{Dipoles, populations, and optical strengths}

The constant and first-order terms of the Cartesian length-gauge
dipole give the Franck--Condon (FC) and Herzberg--Teller (HT) amplitudes,
respectively. For a dipole $\boldsymbol\mu_i$ in $e\,\mathrm{bohr}$ at the
initial origin, define $\mathbf d_k=\partial_{Q_{ik}}\boldsymbol\mu$ there.
The coordinate operator connects the two neighboring initial occupations,
so the complete linear amplitude is
\begin{equation}
 \begin{aligned}
 \boldsymbol\mu_{\mathbf m\mathbf n}={}&
 \boldsymbol\mu_i F_{\mathbf n\mathbf m}
 +\sum_k \frac{\mathbf d_k}{\sqrt{2\omega_{ik}}}
 \left(\sqrt{n_k+1}\, F_{\mathbf n+\mathbf e_k,\mathbf m}
       +\sqrt{n_k}\, F_{\mathbf n-\mathbf e_k,\mathbf m}\right),\\
 S_{\mathbf m\mathbf n}={}&\|\boldsymbol\mu_{\mathbf m\mathbf n}\|^2 .
 \end{aligned}
 \label{eq:strength}
\end{equation}
The FC and HT vectors interfere before the Cartesian components are
squared and summed. Both terms must use the same electronic phase and
Cartesian frame. A common phase leaves $S$ invariant. The relative
phase of the dipole and its derivatives controls the interference.
The derivative rows use the signs and order of the chosen normal modes.
For a single ground-endpoint expansion reused in emission, take
$Q_e=JQ_g+K$ in the absorption convention. With $D_g$ the matrix whose rows
are the ground-endpoint derivatives, the inverse affine map gives
$D_e=JD_g$ and $\boldsymbol\mu_e=\boldsymbol\mu_g-D_g^{\mathsf T}J^{\mathsf T}K$
for that same linear function. Prepared spectra may use independent dipole
expansions evaluated at the two initial endpoints.\cite{baiardi}

The signed photon energy is $E_\gamma=\eta(\Delta_{00}+
\mathbf m\cdot\boldsymbol\omega_f-\mathbf n\cdot\boldsymbol\omega_i)$,
with $\eta=+1$ for absorption and $\eta=-1$ for emission, for which
the excited surface is initial. Optical lines require $E_\gamma>0$.
At temperature $T>0$, the initial occupation carries
$p_{\mathbf n}=\exp(-\beta\,\mathbf n\cdot\boldsymbol\omega_i)
\prod_k(1-e^{-\beta\omega_{ik}})$ with $\beta=(k_{\mathrm B}T)^{-1}$.
This uses the full harmonic partition function. At 0~K only the initial
vacuum is occupied. For population coverage the assigned-stick solver
sums $p_{\mathbf n}$ over retained initial occupations. Final strength
coverage compares the retained final sum $\sum_{\mathbf n\mathbf m}
p_{\mathbf n}S_{\mathbf m\mathbf n}$ with the complete final-state closure
for the same initial states, both before the positive-energy cut. The
closure for each $\mathbf n$ is
$\|\boldsymbol\mu_i\|^2+\sum_k\|\mathbf d_k\|^2(2n_k+1)/(2\omega_{ik})$,
from the initial number-state expectation of $\|\boldsymbol\mu(Q_i)\|^2$.

For $E_\gamma$ in Hartree and $S$ in squared atomic units, the dimensionless
absorption oscillator strength and the vacuum spontaneous emission rate are
\begin{equation}
 f=\frac{2}{3}E_\gamma S,
 \qquad A=\frac{4\alpha^3}{3t_{\mathrm{au}}}E_\gamma^3 S ,
 \label{eq:optical-weight}
\end{equation}
where $\alpha$ is the fine-structure constant and $t_{\mathrm{au}}$ is the
atomic unit of time in seconds. Thermal sticks carry $pf$ or $pA$, with
the population conditional on the initial electronic surface. The
population-weighted Einstein coefficients count photons per second per
molecule in the emitting state. Multiplication by the photon energy in
joules converts photon rate to radiant power. For absorption per unit
wavenumber, the oscillator-strength density gives the cross section on
multiplication by $\pi r_e$, where $r_e$ is the classical electron radius
in cm.

\subsection{Bounds for integrated harmonic observables}

The population and dipole closure tests precede photon weighting.
The corresponding positive-photon moments of the harmonic model are
\begin{equation}
 M_k=\sum_{\mathbf n,\mathbf m}p_{\mathbf n}S_{\mathbf m\mathbf n}
       [E_\gamma]_+^k,\qquad [x]_+=\max(0,x).
 \label{eq:harmonic-photon-moments}
\end{equation}
The absorption oscillator-strength sum is proportional to $M_1$,
the radiative rate to $M_3$, and the radiant power to $M_4$.
The finite stick basis yields $M_k^{\mathrm{ret}}$. With upper bounds
$U_{i,k}$ for contributions from omitted initial states and $U_{f,k}$
for omitted final states reached from retained initial states, the
coverage is bounded below by
\begin{equation}
 \frac{M_k^{\mathrm{ret}}}{M_k^{\mathrm{ret}}+U_{i,k}+U_{f,k}}
 \label{eq:harmonic-moment-coverage}
\end{equation}
for a positive denominator.

For emission, set $E_{\mathbf n}=\mathbf n\cdot\boldsymbol\omega_i$ and
let $C_{\mathbf n}=\langle\|\boldsymbol\mu(Q_i)\|^2\rangle_{\mathbf n}$
be the full dipole closure in initial level $\mathbf n$. An omitted
final level above a total-quanta cutoff $N_f$ has excitation energy
of at least $(N_f+1)\min_j\omega_{fj}$. Thus its positive photon
energy is at most
$[E_{\mathbf n}-\Delta_{00}-(N_f+1)\min_j\omega_{fj}]_+$.
For each retained initial level, the $k$th power of this ceiling
multiplies the unrepresented dipole strength
$C_{\mathbf n}-\sum_{\mathbf m\ {\rm retained}}S_{\mathbf m\mathbf n}$.
The population-weighted sum gives the final-state bound. For omitted
initial levels, $E_{\mathbf n}+[-\Delta_{00}]_+$ bounds the positive
photon energy to every final level. The corresponding bound is the
thermal average of $C_{\mathbf n}(E_{\mathbf n}+[-\Delta_{00}]_+)^k$
over the full oscillator ensemble minus its retained initial terms.
The full average follows from the independent geometric mode
occupations and the affine dependence of $C_{\mathbf n}$ on their
occupation numbers.

For absorption, final-state closure gives the signed first moment
at fixed $\mathbf n$,
\begin{equation}
 \sum_{\mathbf m} E_\gamma S_{\mathbf m\mathbf n}
 = \bigl\langle\|\boldsymbol\mu(Q_i)\|^2 (V_f-V_i)\bigr\rangle_{\mathbf n}
   +\tfrac12\sum_j\|\mathbf d_j\|^2 .
 \label{eq:harmonic-signed-first-moment}
\end{equation}
Here $V_i$ and $V_f$ include the respective electronic minima and
harmonic potentials expressed in the same initial coordinates. The
last term comes from the common kinetic operator acting on the affine
dipole. The signed contributions of each omitted set follow by
subtracting retained sums from the corresponding analytic closure,
using Gaussian coordinate averages for the full thermal ensemble.
The positive first moment equals the signed first moment plus the
magnitude of its negative-energy part. The same energy ceilings as
for emission bound that negative part at $k=1$. Adding those bounds
to the signed contributions gives $U_{i,1}$ and $U_{f,1}$.

The numerical bounds include allowances for roundoff in differences
of totals and retained contributions. The adaptive harmonic stick
solver increases the two quanta cutoffs until the population, dipole
closure when defined, and photon-moment coverages meet the target
(99.5\% by default), or reports the unmet criterion. If the overlap
budget blocks further growth, the solver searches the affordable basis
pairs using the largest feasible final cutoff for each initial cutoff
that can meet the population target. For finite-temperature HT it counts
the extra initial overlap layer needed for coordinate insertion. These bounds
concern integrated observables within the chosen harmonic and dipole
model. Individual band strengths and broadened peak positions require
their own convergence checks.

\subsection{Selected anharmonic coordinates}

For a single vibrational mode the automatic calculation samples both
potentials and the signed vector dipole on a common mass-weighted
coordinate. It solves each Hamiltonian by a sine discrete variable
representation (DVR) with Dirichlet walls. It checks the interpolation
against independent electronic samples and tests the nuclear levels and
line strengths on finer grids. The domain is extended if the retained
levels reach the potential at the walls or retain significant probability
near an edge. These tests distinguish the precision of the interpolated
electronic surfaces from convergence of the nuclear solution.\cite{dvr}

For multiple modes, the automatic calculation uses local cuts about both
minima in dimensionless coordinates $x_k=\sqrt{\omega_k}Q_k$. Each fit
retains the Hessian quadratic term $\sum_k\omega_k x_k^2/2$ and tests cubic
and quartic corrections along individual axes and selected pairs and
triples. Held-out axis, mixed, and collective displacements constrain the
sampled model. Vibrational configuration interaction (VCI) diagonalizes
the correction in a harmonic product basis. Smoothly damping the
corrections at the edge of the sampled region restores confining harmonic
tails. The calculation checks the retained energies and eigenvectors or
degenerate subspaces against larger bases and altered damping radii. The
transition amplitudes contract both VCI eigenvectors with the harmonic
FC and HT matrix elements. Unprobed couplings remain outside the test,
and a verified second minimum requires a different nuclear model.

At positive temperature the automatic calculation weights transitions by
the initial surface's Boltzmann populations. For one-mode DVR, the full
finite-grid eigenvalue partition is checked against finer grids and a
larger domain. The retained thermal population must remain confined below
the sampled well walls. In VCI, initial levels are selected in energy
order. An operator lower bound controls the Boltzmann weight omitted from
the basis of the tapered Hamiltonian, and basis and taper changes test
the stability of the retained populations. The workflow checks both
initial population coverage and final dipole-strength closure before
returning either spectrum.

In the prepared formaldehyde example, the active space includes the
reflection-odd inversion displacement between the two $S_1$ wells, the
reflection-even mean displacement of the wells from the $S_0$ minimum,
and the independent planar relaxation direction toward the $S_1$ saddle.
The RHF--CIS/6-31G(d) potentials are even in the inversion coordinate and
the retained dipole component is odd. A tensor sine DVR solves the coupled
three-coordinate motion with a constant mass-weighted Cartesian kinetic
operator. The three remaining modes contribute only their differential
harmonic zero-point energy. At 0~K the initial even level therefore
couples to odd levels on the other surface, and the even-to-even 0--0
transition is dark within this model. The profile uses a Lanczos spectral
quadrature of the final Hamiltonian with starting vector
$\mu\chi_{i0}$, checked against quadrature order and changes in the nuclear
grid and box. The model retains the double well and planar relaxation
in its active space.

\subsection{Periodic rotor with harmonic spectators}

For a supplied closed angular path, a Fourier DVR uses
$H_s^{\mathrm{rot}}=-(2I_s)^{-1}\partial_\theta^2+V_s(\theta)$ with $\theta$
in radians. The potentials share an electronic energy reference. The
uniform grid spans one full cycle without duplicating the endpoint. For
a path $R_s(\theta)$ on a nonlinear molecule, the local inertia in
$m_e\,\mathrm{bohr}^2$ is $\|P_{\mathrm{int}}M^{1/2}\partial_\theta R_s\|^2$,
where $M$ is the diagonal Cartesian mass matrix and $P_{\mathrm{int}}$
projects off the six mass-weighted rigid motions. The prepared assembly
uses a surface-specific constant $I_s$ with $I_s^{-1}$ equal to the
uniform-angle mean local inverse inertia. The stationary mass-weighted
Hessians restricted to the subspaces perpendicular to the rigid and
torsional directions supply the harmonic spectators. If $\hat t$ is
the normalized projected torsion tangent, this restriction has projector
$P_{\mathrm{sp}}=P_{\mathrm{int}}-\hat t\hat t^{\mathsf T}$.
Let $W_{j+1}$ be an orthonormal basis of that subspace at the next angle
and $\mathcal R_j$ the atom-block rotation of the moving top. For the
transported modes $U_j$ the local update is
\begin{equation}
 U_{j+1}=W_{j+1}\operatorname{polar}
          \bigl(W_{j+1}^{\mathsf T}\mathcal R_j U_j\bigr),
 \label{eq:transport}
\end{equation}
where $\operatorname{polar}(B)$ denotes the orthogonal polar factor of $B$.
The singular values before reorthogonalization quantify local subspace
loss. The bases transported into a common angular frame give the spectator
Duschinsky map and displacement.
The separable Hamiltonian omits angular changes of the spectator
Hamiltonian, variable inertia, and rotor--spectator coupling.
The population model uses unit nuclear-spin statistical weights.\cite{hybriddvr}

For a prepared path with signed overlaps between physical electronic
states, the cumulative signs of successive overlaps define a continuous
gauge for each isolated real state on the sampled path. The product
including the closing link against the original first-point gauge gives
its closure sign $h_s$. This contract requires physical-state overlaps
for both states. The particle--hole similarity in
Eq.~\eqref{eq:electronic-tracking} tracks the transition amplitude and
supplies its phase for the signed dipole.
Single-valuedness of the product of electronic and nuclear factors then
requires $\chi_s(\theta+2\pi)=h_s\chi_s(\theta)$. The Fourier kinetic basis
has integer wave numbers for $h_s=+1$ and half-integer wave numbers for
$h_s=-1$. The potential remains periodic, while the signed dipole has
closure sign $h_i h_f$ in the transported gauge.\cite{longuethiggins}
These signs make the product $\chi_f^*\boldsymbol\mu\chi_i$ periodic.
For DVR eigenvector components $C_{sj,a}$ normalized with the quadrature
weight already included, the active dipole vector is
$\mathbf A_{ba}=\sum_j C_{fj,b}^*\boldsymbol\mu(\theta_j)C_{ij,a}$.
Contributions at different angles are summed coherently, so their
relative signs matter even though the overall gauge sign of a line
does not affect its strength. The rotor-only dipole approximation gives
the total amplitude $\mathbf A_{ba} F_{\mathbf n\mathbf m}$, where
$F_{\mathbf n\mathbf m}$ is now the spectator overlap. A finer DVR grid
on Fourier-interpolated inputs checks the nuclear discretization.
Convergence of the underlying cuts requires independent electronic
samples.

\subsection{All-quanta thermal profiles}

Let $\epsilon_{sa}$ be a rotor eigenenergy and $p^{\mathrm{rot}}_a$ its
initial Boltzmann population, normalized over the full finite DVR spectrum.
For a spectator-independent dipole, the signed thermal dipole measure is
the convolution $\mathcal S_{\mathrm{rot}}*\mathcal S_{\mathrm{h}}$ of
\begin{align}
 \mathcal S_{\mathrm{rot}}(x)&=
 \sum_{ab} p^{\mathrm{rot}}_a\|\mathbf A_{ba}\|^2
  \delta\bigl(x-\eta(\epsilon_{fb}-\epsilon_{ia})\bigr),\\
 \mathcal S_{\mathrm{h}}(x)&=
 \sum_{\mathbf n\mathbf m} p_{\mathbf n}|F_{\mathbf n\mathbf m}|^2
  \delta\bigl(x-\eta[\Delta Z_{\mathrm h}+\mathbf m\cdot\boldsymbol\omega_f
                           -\mathbf n\cdot\boldsymbol\omega_i]\bigr),
\end{align}
where $\Delta Z_{\mathrm h}$ is the spectator zero-point difference. At
0~K the degenerate initial rotor ground manifold, when present, is
uniformly populated. The retained rotor levels determine active population
and strength coverage.
The harmonic measure has unit integral by completeness, includes all
initial and final harmonic occupations, and permits unequal frequencies,
displacement and Duschinsky mixing.

The transform of the harmonic measure is a Gaussian contraction of two
copies of Eq.~\eqref{eq:fc-generating}. Define $L(t)$ as diagonal with
entries $\exp[(-\beta+i\eta t)\omega_{ik}/2]$ for initial components and
$\exp[-i\eta t\omega_{fk}/2]$ for final components, where $t$ is conjugate
to energy in the Fourier transform. Let $B_t=L B L$ and $\ell_t=L\ell$.
Contracting the two copies with equal Fock indices gives
\begin{equation}
 \widehat{\mathcal S}_{\mathrm h}(t)=
 F_{\mathbf 0}^2\prod_k(1-e^{-\beta\omega_{ik}})
 \frac{\exp\!\bigl[-it\eta\Delta Z_{\mathrm h}+
      \ell_t^{\mathsf T}(I-B_t)^{-1}\ell_t\bigr]}
      {\sqrt{\det(I-B_t^2)}} .
 \label{eq:thermal-characteristic}
\end{equation}
The square root is continued from $t=0$. In the 0~K limit the initial
activities vanish. In the standalone harmonic profile, the coordinate
insertion $(z_k+\partial_{z_k})/\sqrt{2\omega_{ik}}$ multiplies $G(z)$ by
an affine vector polynomial, including the constant dipole. The same
Gaussian contraction gives the squared mean of this polynomial plus a
cross-covariance term, with covariance $(I-B_t^2)^{-1}$ for the variables
after scaling by $L(t)$. This sums the FC and linear HT terms coherently
and recovers thermal dipole closure at $t=0$. The spectator-independent
rotor convolution uses the FC contraction in Eq.~\eqref{eq:thermal-characteristic}.
\cite{huhberger,baiardi}

For equal, undisplaced harmonic surfaces the Condon origin and the
linear HT sidebands give the thermal lines directly. The solver selects
positive photon energies, weights them by $E_\gamma$ or $E_\gamma^3$,
and centers the Gaussians at the exact line energies.
In the general Fourier calculation, the active and harmonic transforms
are multiplied before imposing the positive-energy cut and photon
weight so that the cutoff acts on the combined energy. The calculation
uses a Fourier mesh in wavenumber, a two-sided
exponential-moment bound on the underlying dipole measure, and a narrow
auxiliary Gaussian of standard deviation twice the mesh spacing. To weight
the original energies, $(u+\sigma_{\rm aux}^2\partial_u)^p$ acts on the
auxiliary-smoothed density, with $p=1$ or $3$. Here $u$ and $\sigma_{\rm aux}$
are in the same wavenumber units, and the result is converted to Hartree
weighting. The remaining Gaussian variance is the requested physical
variance minus the auxiliary variance. The solver rejects the mesh if
estimated leakage across zero or the negative photon-weighted bins
exceed $10^{-8}$ of the corresponding positive measure. The harmonic
sum is exact within the separable model. The profile retains a mesh error,
especially near zero photon energy.
